\section{An explicit lattice comparison for Gaussian determinant three}
\label{sec:gaussian-nonsplit-lattice}

Put $\mathcal O=\mathbb Z[i]$, $K=\mathbb Q(i)$, and
\[
 H_0=\begin{pmatrix}0&I_2&0\\I_2&0&0\\0&0&J_3\end{pmatrix},
             \qquad J_3=\operatorname{diag}(1,-3).
\]
We use conjugate-linear first arguments.  A Hermitian form of
signature $(3,3)$ has signed determinant $-\det h$ modulo
$\operatorname{Nm}_{K/\mathbb Q}(K^*)$.

\subsection*{D-E8-LAT-01. Two-sided denominators bounded by the polarization volume}

Let $\Lambda$ be a free rank-six $\mathcal O$-lattice with an
integral nondegenerate Hermitian form $h$ of signature $(3,3)$
and signed determinant class three.  Write
$\Delta=|\det h|\in\mathbb Z_{>0}$.  There are an
$\mathcal O$-basis of $\Lambda$, a matrix $T\in\operatorname{GL}_6(K)$,
and an integer $q>0$ such that
\[
 T^*hT=H_0,\qquad qT,\ qT^{-1}\in M_6(\mathcal O),\qquad
                        q\le69984\Delta^3.                  \tag{LAT.1}
\]
Both inclusion directions are part of the assertion.  The number
$69984\Delta^3$ is an upper bound for a clearing integer; it is
not claimed to be itself divisible by every possible clearing
integer.

\begin{proof}
Landherr's published classification identifies the rational
Hermitian space with $H_0$: see \cite[p.~246 for the invariants,
first paragraph of p.~247 for their completeness]{E8LAT:Landherr}.
Only this existence statement is used; no denominator bound is
deduced from it.  In particular the rational Witt index is two.
The binary plane $J_3$ is anisotropic since three is not a
Gaussian norm.

\emph{Step1: split the two hyperbolic planes with controlled
denominators.}  Intersect a rational maximal isotropic two-plane
with $\Lambda$.  The intersection is primitive.  Since
$\mathcal O$ is a PID, choose its basis $v=(v_1,v_2)$ and extend
it to a basis of $\Lambda$.  The pairing with the other four
basis vectors has rank two.  Smith normal form over $\mathcal O$
allows further integral basis changes so that, in blocks
$(v,w,z)$ of size two, the Gram matrix is
\[
 \begin{pmatrix}
 0&B&0\\ B^*&F&E\\0&E^*&A
 \end{pmatrix},\qquad
             B=\operatorname{diag}(b_1,b_2),\quad b_1b_2\ne0.
                                                                  \tag{LAT.2}
\]
All its entries are integral, and $F,A$ are Hermitian.  Put
$b=\det B$, $n_b=b\overline b$, and $d=|\det A|$.  Taking
determinants of (LAT.2) gives
\[
                   \Delta=n_b d.                            \tag{LAT.3}
\]
The following change of rational basis is explicit:
\[
 \begin{split}
 u&=\left(w-v(B^*)^{-1}F/2\right)B^{-1},\\
 z'&=z-v(B^*)^{-1}E.
 \end{split}                                                   \tag{LAT.4}
\]
It gives Gram matrix $\begin{psmallmatrix}0&I_2&0\\I_2&0&0\\0&0&A\end{psmallmatrix}$.
Indeed it kills the $w$-self-pairing and its pairing with $z'$,
and does not alter the $z'$-self-pairing.  The inverse change is
\[
             w=uB+v(B^*)^{-1}F/2,
             \qquad z=z'+v(B^*)^{-1}E.                        \tag{LAT.5}
\]
The denominators of both (LAT.4) and (LAT.5) divide $2n_b$:
the only products of inverse determinants are
$b^{-1}\overline b^{-1}=n_b^{-1}$, and either individual inverse
is also cleared by $n_b$.  Thus this step has a two-sided
clearing integer $2n_b$, not an uncontrolled basis-height
factor.  The residual integral binary form $A$ has signature
$(1,1)$ and signed determinant class three, hence is anisotropic.

\emph{Step2: reduce the binary coefficients using the Gaussian
Euclidean algorithm.}  Choose a nonzero integral vector on
which $|A(x,x)|$ is the smallest possible positive integer.
Anisotropy ensures these values are nonzero.  This vector is
primitive: division by a nonunit common divisor of its two
coordinates would reduce its value by that divisor's norm.
Extend it to an integral basis.  On replacing $A$ temporarily
by $-A$ if necessary, its matrix is
\[
                    A=\begin{pmatrix}a&z\\\overline z&c\end{pmatrix},
                       \qquad a>0.
\]
Shear the second vector by a Gaussian integral multiple of the
first, rounding both coordinates of $z/a$ to nearest integers.
This gives
\[
       |z|^2\le a^2/2,\qquad |c|\ge a.                       \tag{LAT.6}
\]
The second inequality is minimality; $c=0$ is excluded by
anisotropy.  Since $ac-|z|^2=-d<0$, the case $c>0$ would give
$|z|^2=ac+d>a^2$, contradicting (LAT.6).  Therefore $c<0$ and
\[
 d=a|c|+|z|^2,\qquad a\le\sqrt d,\quad |c|\le d,
 \qquad |\Re z|+|\Im z|\le a\le\sqrt d.                      \tag{LAT.7}
\]
These integral basis changes and their inverses introduce no
new denominators.

\emph{Step3: a complete small-zero argument in five variables.}
Consider the integral quadratic form
\[
 Q(X_1,X_2,X_3,X_4,t)
       =A(X_1+iX_2,X_3+iX_4)-t^2.
\]
Write $Q(x)=x^t\mathcal Qx$ with an integral symmetric matrix,
and put $S=\sum_{j,k}|\mathcal Q_{jk}|$.  Formula (LAT.7) gives
\[
 S=2a+2|c|+4(|\Re z|+|\Im z|)+1
                    \le2d+6\sqrt d+1\le9d.                  \tag{LAT.8}
\]
The form $Q$ is isotropic: by Landherr's classification $A$ is
isometric to $J_3$, which represents one.  It has rational Witt
index exactly one.  If it had a totally isotropic two-plane,
that plane would meet the hyperplane $t=0$ nontrivially,
contradicting the anisotropy of $A$.

Here is the full elementary height argument, so no unverified
version of a small-zero constant is needed.  For an integral
isotropic quadratic form in five variables of Witt index one,
choose a nonzero integral zero $x$ with minimal sup-norm
$M=\max_j|x_j|$.  It is primitive.  Reorder coordinates and
change its sign so that $x_5=M$.  Set
\[
                       Q_0=\lfloor\sqrt{3S}\rfloor+1.
\]
Suppose $M>Q_0^4$.  Among the $Q_0^4+1$ points
\[
 \left(\{j x_1/M\},\ldots,\{j x_4/M\}\right),
                         \quad 0\le j\le Q_0^4,
\]
two lie in the same one of the $Q_0^4$ half-open boxes of
side $1/Q_0$.  Subtract their indices to obtain
$1\le k\le Q_0^4<M$, and subtract their integer parts to
obtain $y\in\mathbb Z^5$ satisfying
\[
 y_5=k,\qquad y=(k/M)x+e,\qquad e_5=0,
                         |e_j|<1/Q_0.
\]
The vectors $x,y$ are not proportional: since $x$ is primitive,
an integral rational multiple of it has integral multiplier,
whereas $0<y_5/x_5=k/M<1$.
Let $\mathcal B$ denote the bilinear form with matrix
$\mathcal Q$.  The reflection numerator
\[
             z_0=Q(y)x-2\mathcal B(x,y)y
                 =Q(e)x-2\mathcal B(x,e)e                     \tag{LAT.9}
\]
is an integral zero of $Q$.  It is nonzero.  If it vanished,
either $y$ would be proportional to $x$, or
$Q(y)=\mathcal B(x,y)=0$; the latter would make their span
a totally isotropic two-plane.  Both alternatives have just
been excluded.  But
\[
 \|z_0\|_\infty
       \le\frac{3SM}{Q_0^2}<M,
\]
contradicting the choice of $x$.  It follows that
\[
 M\le Q_0^4\le(2\sqrt{3S})^4=144S^2.
\]
Applied to our $Q$, this gives a nonzero zero with
\[
              0<|t|\le11664d^2.                             \tag{LAT.10}
\]
The inequality $t\ne0$ follows from the anisotropy of $A$.
The selected maximal coordinate in the height argument need
not be this $t$ coordinate; (LAT.10) follows because every
coordinate is bounded by $M$.

\emph{Step4: an explicit binary isometry and its inverse.}
Put $p=(X_1+iX_2,X_3+iX_4)^t$ and $v=p/t$, so that
$v^*Av=1$.  Write $(\ell_1,\ell_2)=v^*A$ and
$u=(-\ell_2,\ell_1)^t$.  Direct multiplication gives
\[
          v^*Au=0,\qquad u^*Au=\det A=-d,
                    \det(v,u)=1.                             \tag{LAT.11}
\]
The coordinates of both $v$ and $u$ have denominators dividing
$|t|$.

Because $d/3$ is a rational Gaussian norm, the positive integer
$3d$ is a sum of two integral squares.  To justify the integral
assertion, factor $3d$ in $\mathbb Z$: at every prime
$p\equiv3\pmod4$ its valuation is even, as follows from it being
a rational norm.  The Gaussian Euclidean factorization then
constructs $\beta\in\mathcal O$ with
$\beta\overline\beta=3d$.  Set $\alpha=\beta/3$.  The matrix
\[
                  T_A=(v,u/\alpha)
\]
satisfies $T_A^*AT_A=J_3$.  Its entries are cleared by
$d|t|$, since $\alpha^{-1}=\overline\beta/d$.
Its determinant is $\alpha^{-1}$ by (LAT.11), and hence
$T_A^{-1}=\alpha\operatorname{adj}(T_A)$.  The entries in
the first row of this inverse have denominators dividing
$|t|$, and those in the second row have denominators dividing
$3|t|$.  In particular the single integer
\[
                            3d|t|                            \tag{LAT.12}
\]
clears both $T_A$ and $T_A^{-1}$.

If Step2 used $-A$, its construction instead gives Gram
$-J_3$ for the original $A$.  The integral determinant-one
matrix
\[
        U=\begin{pmatrix}1+i&-3\\1&-1+i\end{pmatrix}
\]
satisfies $U^*(-J_3)U=J_3$.  Multiplying by $U$ fixes the
sign without changing either denominator bound.

Combine the binary change with (LAT.4).  Both the total
change and its inverse are cleared by
\[
          q=(2n_b)(3d|t|)=6\Delta|t|
                 \le69984\Delta d^2\le69984\Delta^3.
\]
This proves (LAT.1), including its two inclusion directions.
\end{proof}

\subsection*{D-E8-LAT-02. Fixed-lattice isogeny and cycle-degree transport}

For two such lattices of volumes $\Delta_S,\Delta_T$ there is
a rational Hermitian isometry $P$ from the first to the second
and an integer $N$ such that
\[
 NP\Lambda_S\subset\Lambda_T,\quad
 NP^{-1}\Lambda_T\subset\Lambda_S,\qquad
          N\le69984^2(\Delta_S\Delta_T)^3.                    \tag{LAT.13}
\]
In particular, if the source lattice has volume at most $M^6$,
then
\[
                       N\le69984^2M^{18}\Delta_T^3.           \tag{LAT.14}
\]
For the $E_8$ source one may take conservatively
$M=2|G|=2786918400$, the primitive-projector denominator.

\begin{proof}
Choose $T_S,T_T$ and clearing integers $q_S,q_T$ from
D-E8-LAT-01, and take
$P=T_TT_S^{-1}$ and $N=q_Sq_T$.  Then
$P^*h_TP=h_S$.  Since both $q_TT_T$ and $q_ST_S^{-1}$
are integral, $NP$ is integral.  The other two integral
matrices give the same assertion for $NP^{-1}$.  This proves
(LAT.13); substituting $\Delta_S\le M^6$ gives (LAT.14).

For a compatible integral alternating polarization $E$ on an
$\mathcal O$-stable lattice use
\[
                 h(x,y)=E(x,iy)+iE(x,y).                     \tag{LAT.15}
\]
It is an integral Hermitian form, and $E$ is its coefficient
of $i$.  The determinant of the alternating matrix is
$(\det h)^2$, so $|\det h|$ is precisely the product of the
polarization elementary divisors.  Therefore a rational
Hermitian isometry is a rational symplectic isometry for the
specified polarizations, with no unspecified positive scalar.

For completeness, the stated source bound uses the actual
primitive projector, not just the group order as a heuristic.
Let $\chi$ be the degree-eight character of the natural
representation and put
\[
 \epsilon=\frac8{|G|}\sum_{g\in G}
       (\chi(g^{-1})+\overline{\chi(g^{-1})})g,\qquad
                   e=\epsilon(1-s)/2,
\]
where $s$ is a reflection.  The character trace in parentheses
is integral, so $(2|G|)e\in\mathbb Z[G]$.  This is a primitive
self-adjoint idempotent in the rational $M_8(K)$ factor.
If $L=H_1(J(C),\mathbb Z)$ is the unimodular Jacobian lattice,
the lattice of its image is $\Lambda_S=eL_{\mathbb Q}\cap L$.
Its polarization dual is $eL$: an integral functional on the
primitive sublattice $\Lambda_S$ extends to $L$, and
unimodularity followed by the self-adjoint projection gives
the representing vector in $eL$.  Therefore
$M\Lambda_S^\#\subset\Lambda_S$ for $M=2|G|$.
Each of its six polarization elementary divisors divides $M$,
proving $\Delta_S\le M^6$.  The central group element $iI$
preserves the image abelian variety and its integral lattice,
so this source lattice is indeed $\mathbb Z[i]$-stable.

Fix the two lattices and identify their period domains by $P$.
At corresponding period points, $NP$ is a $K$-linear isogeny
$\phi:B_S\to B_T$ and
\[
                \phi^*\theta_T=N^2\theta_S,
 \qquad \deg\phi=N^{12}\Delta_S/\Delta_T.                    \tag{LAT.16}
\]
The integer $N$ and the rational isometry can be fixed on a
marked polydisc, so this is a family comparison, not an
isogeny chosen with an uncontrolled denominator at each fiber.
For an effective codimension-three source cycle $Z$,
the projection formula gives
\[
                  \deg_{\theta_T}(\phi_*Z)
                      =N^6\deg_{\theta_S}(Z).                \tag{LAT.17}
\]
Thus the numerical upper bound in (LAT.14) may safely be used
to bound transported cycle degrees, even though that upper
bound is not itself claimed to clear every isometry.  Weil
projection is preserved by a rational $K$-isogeny, so a
nonzero detecting projection remains nonzero.
\end{proof}

For a target lattice stable only under the fixed order
$\mathbb Z+f\mathbb Z[i]$, replace it first by
$\Lambda'=f\mathcal O\Lambda\subset\Lambda$.  As proved in
D-G34-07b, $f\Lambda\subset\Lambda'\subset\Lambda$ and
$[\Lambda:\Lambda']\le f^{12}$ in dimension six.  Its pulled-back
polarization has volume at most $f^{12}\Delta_T$, so (LAT.14)
becomes
\[
                   N\le69984^2M^{18}f^{36}\Delta_T^3.         \tag{LAT.18}
\]
The subsequent identity isogeny to the original target pulls
back its polarization exactly, and adds no factor to (LAT.17).
The scope is a fixed conductor $f$, not a bound independent
of all orders.
